\hspace{3mm}

\centerline{\bf \Huge База по графам}

Теория графов - это важная область дискретной математики, изучающая объекты, называемые графами, которые состоят из точек (вершин) и рёбер (связей между вершинами). Графы широко применяются в различных областях, таких как компьютерные науки, логистика, социология, транспорт, биоинформатика и многие другие. 

Начнем с определения графа.

\textbf{Граф} это множество состоящее из точек и линий между ними, которые называются \textit{вершинами} и \textit{ребрами}. Вершины это какие-то объекты, а ребра это связи, взаимоотношения между этими объектами.

Типичные примеры графа это дороги, проведенные между городами или же дружба между людьми $-$ между двумя людьми проводится ребро, если они дружат.

\begin{figure}[h]
    \centering
    \includegraphics[width=0.3\linewidth]{Теория/Graph_picture.png}
    \caption{Пример графа}
\end{figure}

Опираясь на этот рисунок дадим несколько определений, связанных с графами.

\begin{itemize}
    \item \textbf{Степенью вершины} называется количество ребер, выходящих из этой вершины. На рисунки у каждой вершины подписана её степень.
    \item \textbf{Изолированная вершина} это вершина степени 0.
    \item \textbf{Кратные рёбра} это несколько ребер соединяющих одну и ту же пару вершин.
    \item \textbf{Петля} это ребро у которого конечная вершина совпадает с начальной(на рисунке номер 5).
    \item \textbf{Путь} это последовательность ребер, где начало следующего ребра это конец предыдущего.
    \item \textbf{Простой путь} это путь, который не приходит дважды в одну и ту же вершину.
    \item \textbf{Цикл} это путь, начальная и конечная вершина которого совпадают. 
    \item \textbf{Простой цикл} это цикл, который не приходит дважды в одну и ту же вершину.
    \item  Граф называется \textbf{связным}, если из любой его вершин по ребрам можно добраться до любой другой.
\end{itemize}

\textbf{Лемма о \sout{рукож} рукопожатиях.} Сумма степеней всех вершин графа - четное число, равное удвоенному числу его ребер.

Доказательство:

Возьмем пустой граф. Сумма степеней вершин такого графа равна нулю. При добавлении ребра, связывающего любые две вершины, сумма всех степеней увеличивается на 2. Таким образом, сумма всех степеней вершин чётна и равна удвоенному числу рёбер.

\textbf{Великорусское подвешивание графа за вершину.} Назовём подвешиванием (ранжированием) графа за вершину следующую процедуру. Выберем любую вершину графа. Назовём её нулевым уровнем (рангом). На первом уровне расположим все вершины, соединённые с вершиной нулевого уровня. На каждом следующем уровне будем располагать все вершины, соединённые с вершинами предыдущего уровня и не расположенные на более ранних уровнях. Процесс продолжается, пока все вершины не расположатся на уровнях.

Такой приём особенно выгодный, если вы знаете:

1) Что в графе степень каждой вершины больше либо равно(меньше либо равна) какого-то фиксированного числа

2) Что-то про расстояние между любыми двумя вершинами. Например, если вы знаете, что от любой вершины можно добраться любой другой используя не больше чем 2 ребра. Тогда, подвешивая граф за вершину, вы точно можете сказать, что у вас не будет четвертого уровня, иначе расстояние от корневой вершины, до вершины 4 уровня будет уже 3. Если говорить по-умному, то если диаметр графа равен $d$, то в любом подвешивании графа за вершину не более чем $d+1$ уровень.

3) Если мы знаем что-то про циклы, например их полное отсутствие или же отсутствие только каких-то конкретных.

Если в графе вообще нет циклов, то это почти определение дерева, не хватает только связности. Если же граф не связный, но в нем нет циклов, то это просто объединение нескольких деревьев, которое называют \textit{лесом}.

Если в графе, например, нет циклов длины 3 и 4, то внутри второго уровня ребра отсутствуют, иначе бы появился цикл длины 3 и так же не может быть такого, что из второго уровня ребра двух вершин попадают в одну на третьем уровне, иначе бы получился цикл длины 4. Лучше нарисовать себе эту картинку на тетради, чтобы лучше разобраться.

\textit{Остовное дерево} графа это дерево, которое является подграфом данного графа, с тем же числом вершин, что и у исходного графа.

\textbf{Важная теорема.} В любом графе можно выделить остовное дерево.

Доказательство:

\hspace{3mm} \problem{Шаг 1} Подвесим граф за любую вершину

\hspace{3mm} \problem{Шаг 2} Внутри каждого уровня почистим всё от ребёр 

\hspace{3mm} \problem{Шаг 3} Из каждой вершины, начиная со второго уровня наверх (то есть на уровень выше) выходит хотя бы одно ребро. Уберём из них все, кроме одного.

Нетрудно проверить, что полученный граф не будем иметь циклов и будет связным, то есть будет деревом.

Зачастую задачи с подвешиванием удобно решать, используя остовные деревья, чтобы лишние ребра не мешались, если нам интересны только лишь вершины. Но, конечно, больший интерес они вызывают у программистов. Многие алгоритмы обхода графа позволяют построить остовное дерево. 